% TEMPLATE-MILC-v3.tex - plantilla maestra para todas las guias MILC v3
% Ver scripts/generadores/build-guia-milc.py para la lista completa de
% claves esperadas. El script valida que no queden placeholders sin
% reemplazar antes de escribir el archivo final.

\documentclass[11pt,letterpaper]{article}
\usepackage[margin=1.75cm]{geometry}
\usepackage[table]{xcolor}
\usepackage{fontspec}
\usepackage[spanish]{babel}
\usepackage{tikz}
\usepackage[most]{tcolorbox}
\usepackage{tabularx}
\usepackage{array}
\usepackage{fancyhdr}
\usepackage{titlesec}
\usepackage{ragged2e}
\usepackage{enumitem}
\usepackage{microtype}

\setmainfont{Helvetica Neue}
\setsansfont{Helvetica Neue}
\setlength{\parindent}{0pt}
\setlength{\parskip}{6pt}
\hyphenpenalty=200
\exhyphenpenalty=200
\pretolerance=400
\tolerance=2000
\setlength{\emergencystretch}{1em}
\renewcommand{\arraystretch}{1.25}
\setlist{leftmargin=5mm,itemsep=1mm,topsep=1mm}
\newcolumntype{Y}{>{\RaggedRight\arraybackslash}X}

\definecolor{milcMagenta}{HTML}{E6007E}
\definecolor{milcVino}{HTML}{5A0038}
\definecolor{milcTurquesa}{HTML}{0093A5}
\definecolor{milcVerde}{HTML}{58A923}
\definecolor{milcMostaza}{HTML}{E5B400}
\definecolor{milcOcre}{HTML}{B86600}
\definecolor{milcNegro}{HTML}{241816}
\definecolor{milcGris}{HTML}{F7F7F4}
\definecolor{milcLinea}{HTML}{E8E2DD}
\definecolor{milcGradePrimary}{HTML}{84CC16}
\definecolor{milcGradeDark}{HTML}{4D7C0F}
\definecolor{milcGradeSoft}{HTML}{ECFCCB}
\definecolor{milcGradeAccent}{HTML}{CFE7FF}
\definecolor{milcGradeText}{HTML}{FFFFFF}
\color{milcNegro}

\pagestyle{fancy}
\fancyhf{}
\lhead{\small Tecnología e Informática · Bebras}
\rhead{\small Momento 5 · MILC v3}
\cfoot{\small\thepage}
\renewcommand{\headrulewidth}{0pt}

\newtcolorbox{titlebox}[1]{
  enhanced,colback=#1,colframe=#1,arc=7mm,boxrule=0pt,
  left=7mm,right=7mm,top=3mm,bottom=3mm,
  before skip=8pt,after skip=8pt,
  fontupper=\sffamily\bfseries\Large\color{white}
}

\newtcolorbox{softbox}[3]{
  enhanced,colback=#2,colframe=#1,borderline west={4pt}{0pt}{#1},
  arc=2mm,boxrule=.4pt,left=4mm,right=4mm,top=3mm,bottom=3mm,
  before skip=6pt,after skip=8pt,title={#3},coltitle=white,
  colbacktitle=#1,fonttitle=\sffamily\bfseries,fontupper=\RaggedRight,
  attach boxed title to top left={xshift=3mm,yshift=-2mm},
  boxed title style={arc=2mm, boxrule=0pt}
}

\newtcolorbox{drawbox}[2]{
  enhanced,colback=white,colframe=#1,arc=4mm,boxrule=1pt,
  left=5mm,right=5mm,top=4mm,bottom=4mm,before skip=8pt,after skip=8pt,
  title={#2},coltitle=white,colbacktitle=#1,fonttitle=\sffamily\bfseries,
  fontupper=\RaggedRight,
  attach boxed title to top left={xshift=4mm,yshift=-2mm},
  boxed title style={arc=3mm, boxrule=0pt}
}

\newtcolorbox{infoband}[1]{
  enhanced,colback=#1!8,colframe=#1!50,arc=2mm,boxrule=.5pt,
  left=4mm,right=4mm,top=2mm,bottom=2mm,before skip=4pt,after skip=4pt,
  fontupper=\small\RaggedRight
}

% Puente narrativo entre secciones (contrato editorial Conéctate).
% Da continuidad: "Acabas de X. Ahora vas a Y porque Z."
% Visualmente: banda izquierda en color del grado, texto en itálica pequeña.
\newtcolorbox{puentebox}{
  enhanced,colback=milcGradeSoft,colframe=milcGradeDark,
  borderline west={3pt}{0pt}{milcGradeDark},
  arc=0mm,boxrule=0pt,left=5mm,right=4mm,top=2.5mm,bottom=2.5mm,
  before skip=4pt,after skip=8pt,
  fontupper=\itshape\small\color{milcNegro}\RaggedRight
}
\newcommand{\puente}[1]{%
  \begin{puentebox}%
    \textcolor{milcGradeDark}{$\rightarrow$}\hspace{2mm}#1%
  \end{puentebox}%
}

\newcommand{\linead}[1]{\noindent\color{milcLinea}\rule{#1}{0.5pt}\color{milcNegro}}
\newcommand{\respuesta}[1]{\par\vspace{2mm}\foreach \n in {1,...,#1}{\noindent\color{milcLinea}\rule{\linewidth}{0.5pt}\par\vspace{5mm}}\color{milcNegro}}
\newcommand{\checkbox}{\textcolor{milcGradeDark}{\fbox{\phantom{x}}}\hspace{5pt}}

\newcommand{\phasecard}[4]{
\begin{tcolorbox}[enhanced,colback=#3,colframe=#2,arc=3mm,boxrule=.5pt,height=3.05cm,valign=center,halign=center,left=1.5mm,right=1.5mm,top=2mm,bottom=2mm]
{\sffamily\bfseries\fontsize{22}{24}\selectfont\textcolor{#2}{#1}}\par
{\sffamily\bfseries\small #4}
\end{tcolorbox}
}

\begin{document}
\thispagestyle{empty}
\begin{tikzpicture}[remember picture,overlay]
  \fill[white] (current page.south west) rectangle (current page.north east);
  \path[fill=milcGradePrimary,rounded corners=16mm]
    ([xshift=.55cm,yshift=-.55cm]current page.north west) rectangle
    ([xshift=-.55cm,yshift=3.65cm]current page.south east);

  \node[anchor=north west,text=milcGradeText] at ([xshift=2.0cm,yshift=-1.85cm]current page.north west)
    {\sffamily\bfseries\fontsize{14}{16}\selectfont MOMENTO};
  \node[anchor=north west,text=milcGradeText,opacity=.82] at ([xshift=2.0cm,yshift=-2.75cm]current page.north west)
    {\sffamily\bfseries\fontsize{9}{11}\selectfont BEBRAS · PENSAR COMO UNA MÁQUINA};

  \begin{scope}[shift={([xshift=-3.05cm,yshift=-2.55cm]current page.north east)}]
    \fill[white,opacity=.80] (-.62,-.52) rectangle (.62,.52);
    \draw[milcGradeText,line width=1pt] (-.62,-.52) rectangle (.62,.52);
    \fill[milcGradeDark] (-.42,-.38) rectangle (-.22,.06);
    \fill[milcGradeAccent] (-.10,-.38) rectangle (.10,.24);
    \fill[milcGradePrimary!60!black] (.22,-.38) rectangle (.42,.38);
  \end{scope}

  \node[anchor=west,text=milcGradeText] at ([xshift=1.9cm,yshift=-12.85cm]current page.north west)
    {\sffamily\bfseries\fontsize{124}{124}\selectfont 5};
  \draw[milcGradeText,line width=1.7pt]
    ([xshift=4.52cm,yshift=-11.68cm]current page.north west) circle (.13cm);

  \node[anchor=west,text=milcGradeText,text width=8.7cm,align=left] at ([xshift=10.35cm,yshift=-11.6cm]current page.north west)
    {
     {\sffamily\bfseries\fontsize{19}{22}\selectfont El idioma de los datos --- binario,\\listas, tablas y árboles}\\[4mm]
     {\sffamily\bfseries\fontsize{23}{26}\selectfont Bebras · Momento 5}\\[2mm]
     {\sffamily\fontsize{13}{16}\selectfont Curso «Pensar como una máquina» · Pensamiento computacional}\\[5mm]
     {\sffamily\bfseries\fontsize{11}{13}\selectfont Institución Educativa Sor María Juliana}\\[1mm]
     {\sffamily\fontsize{10}{12}\selectfont Autor: Dr. Álvaro Cárdenas Orozco}
    };

  \path[fill=milcGradeSoft,rounded corners=7mm]
    ([xshift=1.35cm,yshift=.85cm]current page.south west) rectangle
    ([xshift=-1.35cm,yshift=4.25cm]current page.south east);
  \draw[milcGradeDark,line width=.65pt,rounded corners=7mm]
    ([xshift=1.35cm,yshift=.85cm]current page.south west) rectangle
    ([xshift=-1.35cm,yshift=4.25cm]current page.south east);
  \node[anchor=center,text=milcNegro,text width=17.0cm,align=center] at ([yshift=2.55cm]current page.south)
    {\sffamily\fontsize{10}{12}\selectfont Metodología MILC · Apertura ancestral · Pensamiento computacional · Triángulo Dussel-Estoicismo-Floridi\\
    \textcolor{milcGradeDark}{\bfseries Producto:} Tu párrafo «Lo mío en dos estados» (EXPLICA) y tu tabla «Mi código de tambor» con 4 mensajes codificados en golpes y silencios (APLICA).};
\end{tikzpicture}
\mbox{}
\newpage

\section*{Apertura: El idioma de los datos: binario, listas, tablas y árboles}

\begin{softbox}{milcGradeDark}{milcGradeSoft}{Saber ancestral}
En los ríos del Pacífico colombiano, mucho antes del celular, las comunidades \textbf{afrocolombianas} se hablaban a distancia con \textbf{tambores}. Un tamborero experto no enviaba palabras: enviaba \emph{patrones}. Combinando golpes y silencios --- un toque fuerte y otro suave, uno largo y otro corto --- armaba avisos que el río llevaba de vereda en vereda: «llegó la canoa», «viene la creciente», «esta noche hay fiesta». El secreto es asombroso por lo simple: con apenas \textbf{dos estados} --- suena o no suena --- se podía decir casi cualquier cosa, siempre que emisor y receptor acordaran antes qué significaba cada combinación. No hacía falta un alfabeto entero; bastaban dos señales bien ordenadas. Ese es, sin que nadie lo bautizara así, el mismo principio con el que hoy una máquina guarda una foto, una canción o tu nombre: \textbf{todo cabe en dos estados}.
\end{softbox}

\begin{softbox}{milcTurquesa}{milcGris}{Saber contemporáneo}
\textbf{Representar datos} es decidir \emph{cómo} escribimos la información para guardarla o enviarla. La gran idea es sorprendente: \textbf{con pocos símbolos se puede representar todo}. El tamborero del Pacífico lo probó con apenas dos --- golpe y silencio ---. El computador hace lo mismo: guarda fotos, audios y textos con solo dos estados, \textbf{encendido (1) y apagado (0)}. Lo importante no es la cantidad de símbolos, sino el acuerdo de qué significa cada combinación. Además del binario, los datos se organizan en \textbf{listas} (orden que importa), \textbf{tablas} (filas y columnas que se cruzan) y \textbf{árboles} (ramas que se abren). Estas cuatro formas aparecen constantemente en los retos Bebras; entenderlas es la llave para resolverlos.
\end{softbox}

\begin{softbox}{milcMostaza}{milcGris}{Pregunta-puente}
\textit{Si el tambor podía decir tanto con solo dos sonidos --- golpe y silencio ---, ¿será que un computador guarda TODO lo que ves --- fotos, audios, mensajes, videos --- usando también apenas dos estados?}
\end{softbox}

\begin{softbox}{milcVerde}{milcGris}{Saber y saber hacer}
Al terminar podrás: (1) \textbf{identificar} cuatro formas de representar datos --- binario, listas, tablas y árboles --- y dar un ejemplo propio de cada una; (2) \textbf{explicar} con tus palabras cómo la información se reduce a dos estados (0 y 1), conectando la idea con el tambor del Pacífico; (3) \textbf{leer} el valor de un número binario de tres posiciones sumando lo encendido; (4) \textbf{aplicar} la idea de codificar con dos estados creando tu propio código de tambor para cuatro mensajes.
\end{softbox}

\small
\begin{tabularx}{\textwidth}{>{\bfseries}p{3.0cm}Y}
\rowcolor{milcGradePrimary!12}Autor & Dr. Álvaro Cárdenas Orozco\\
DBA & Desarrolla el pensamiento computacional --- descomposición, patrones, abstracción y algoritmos --- para resolver problemas (transversal 6°--11°, alineado a los lineamientos MEN de Tecnología e Informática)\\
Referentes & Valentina Dagienė (Bebras) · Pensamiento computacional (Wing) · Dussel · Estoicismo · Floridi · Saberes del Valle y el Pacífico\\
Producto final & Tu párrafo «Lo mío en dos estados» (EXPLICA) y tu tabla «Mi código de tambor» con 4 mensajes codificados en golpes y silencios (APLICA).\\
\end{tabularx}
\normalsize

\puente{Ya sabes que el tambor del Pacífico codifica mensajes con dos estados. Ahora mira el mapa de la sesión: vas a ir de ese saber a entender el binario y las cuatro formas en que los datos se organizan.}

\begin{titlebox}{milcGradeDark}
Ruta MILC: así vamos a aprender
\end{titlebox}
\begin{tabularx}{\textwidth}{>{\centering\arraybackslash}X>{\centering\arraybackslash}X>{\centering\arraybackslash}X>{\centering\arraybackslash}X}
\phasecard{1}{milcVerde}{milcGris}{Escucha\\[-1mm]\small Observo, escucho y pregunto.} &
\phasecard{2}{milcTurquesa}{milcGris}{Sistematización\\[-1mm]\small Pensamiento computacional.} &
\phasecard{3}{milcMagenta}{milcGris}{Praxis\\[-1mm]\small Construyo y aplico.} &
\phasecard{4}{milcVino}{milcGris}{Evaluación\\[-1mm]\small Reflexión liberadora.}
\end{tabularx}


\puente{Antes de nombrar las formas de representar datos, conviene verlas en acción. Empieza por los toques de tambor: ¿cuántas cosas distintas se pueden decir con solo dos señales?}

\begin{titlebox}{milcVerde}
Fase 1 · Escucha
\end{titlebox}

\begin{softbox}{milcVerde}{milcGris}{Escena para escuchar}
\textbf{Imagina a un tamborero del Pacífico} en la orilla del río. Solo puede hacer dos cosas: golpear (G) o callar (S, de silencio). Tiene que enviar al menos tres avisos distintos a otra vereda usando únicamente esas dos señales en combinación. Fíjate: ¿cuántas combinaciones de dos posiciones puedes armar con G y S? (GG, GS, SG, SS) ¿Y de tres posiciones? Antes de leer más, tómate un momento y anótalas.
\end{softbox}

\checkbox Encontraste al menos cuatro combinaciones distintas de dos posiciones (GG, GS, SG, SS).\\
\checkbox Señalaste que cambiar el orden cambia el significado: GS no es lo mismo que SG.\\
\checkbox Notaste que al agregar una posición más, el número de combinaciones se duplica.

\begin{infoband}{milcVerde}
\textbf{En tu cuaderno (Actividad 1 · IDENTIFICA).} Título: «Actividad 1 --- Los toques del río». \textbf{Formato:} lista de combinaciones + una oración de conclusión. \textbf{Extensión:} las combinaciones de dos posiciones y al menos una observación sobre qué pasa al agregar una tercera. \textbf{Sabes que terminaste cuando} tu lista muestra todas las combinaciones de dos posiciones (son cuatro) y tu oración dice, con tus palabras, por qué el orden importa.
\end{infoband}

\puente{Acabas de explorar la idea de codificar con dos estados. Ahora vamos a nombrar y entender las cuatro representaciones una a una, porque con ellas vas a resolver los retos que siguen.}

\begin{titlebox}{milcTurquesa}
Fase 2 · Sistematización con pensamiento computacional
\end{titlebox}

\begin{softbox}{milcTurquesa}{milcGris}{Cuatro pilares aplicados al tema}
Los datos no flotan sueltos: siempre están \textbf{representados}, es decir, escritos en alguna forma acordada para poder guardarlos o enviarlos. Con pocas formas básicas se organiza \emph{toda} la información de un computador. Aquí están las cuatro que más vas a encontrar en Bebras, con ejemplos de tu entorno en Cartago.
\end{softbox}

\small
\begin{tabularx}{\textwidth}{>{\bfseries\RaggedRight\arraybackslash}p{3.6cm}Y}
\rowcolor{milcTurquesa!90}\color{white}Pilar & \color{white}\bfseries Aplicación al tema\\
1. Descomposición & \textbf{Binario (0 y 1):} solo dos símbolos, como el tambor que suena o calla. Cada posición vale el \textbf{doble} de la de su derecha (las tres primeras: 4, 2, 1). Se suman las posiciones encendidas. Ejemplo: 101 = 4 + 0 + 1 = 5. Con tres posiciones se representan los números del 0 al 7.\\
2. Reconocimiento de patrones & \textbf{Listas:} datos puestos en un orden que importa. La fila para el bus en el terminal de Cartago es una lista: el primero llega primero. Cambiar el orden cambia el significado.\\
3. Abstracción & \textbf{Tablas:} datos organizados en filas y columnas para cruzarlos. El horario de tu salón es una tabla: una casilla te dice qué clase hay el martes a la segunda hora.\\
4. Algoritmo conceptual & \textbf{Árboles:} datos que se ramifican, donde cada nodo abre otros. El árbol genealógico de tu familia es un árbol: de tus abuelos salen tus papás, y de ellos sales tú y tus hermanos.\\
\end{tabularx}
\normalsize

\begin{drawbox}{milcTurquesa}{Cómo leer un número binario de tres posiciones}
\textbf{1. Escribe las posiciones} de derecha a izquierda: 1, 2, 4. \\[1mm] \textbf{2. Anota el dígito} (0 o 1) sobre cada posición. \\[1mm] \textbf{3. Multiplica} cada dígito por su posición: encendido suma su valor, apagado suma cero. \\[1mm] \textbf{4. Suma} los resultados. Ejemplo: 110 encendido en 4 y en 2, apagado en 1: 4 + 2 + 0 = \textbf{6}. \\[1mm] \textbf{5. Comprueba} que el resultado está entre 0 y 7 (rango de 3 bits).
\end{drawbox}

\begin{softbox}{milcMostaza}{milcGris}{Errores comunes y su corrección}
\textbf{Leer el binario de izquierda a derecha sin darle valor a cada posición:} la posición más a la izquierda vale 4, no 1. \textbf{Confundir listas con tablas:} una lista tiene un solo orden; una tabla cruza dos dimensiones. \textbf{Creer que los árboles son solo para familia:} cualquier dato que se ramifique --- categorías, menús, carpetas --- es un árbol.
\end{softbox}

\begin{infoband}{milcTurquesa}
\textbf{En tu cuaderno (Actividad 2 · EXPLICA).} Título: «Actividad 2 --- Las cuatro formas con mis palabras». \textbf{Formato:} tabla de 2 columnas (representación · ejemplo de mi día). \textbf{Extensión:} las cuatro representaciones, cada una con un ejemplo distinto a los del texto. \textbf{Sabes que terminaste cuando} un compañero entendería tus ejemplos sin que se los expliques.
\end{infoband}

\puente{Ya distingues binario, listas, tablas y árboles. Es hora de usarlos: vas a codificar algo tuyo con dos estados y luego inventar tu propio código de tambor.}

\begin{titlebox}{milcMagenta}
Fase 3 · Praxis con pensamiento computacional
\end{titlebox}

\begin{softbox}{milcMagenta}{milcGris}{Construyendo el producto con los cuatro pilares}
Las representaciones no son adorno: se usan. Vas a resolver un \textbf{reto estilo Bebras sobre binario}, aplicando los cuatro pasos para leer un número de tres posiciones. Lee con calma --- la clave está en el valor de cada posición, no en adivinar.
\end{softbox}

\small
\begin{tabularx}{\textwidth}{>{\bfseries\RaggedRight\arraybackslash}p{3.6cm}Y}
\rowcolor{milcMagenta!90}\color{white}Pilar & \color{white}\bfseries Aplicación al producto\\
Descomposición & \textbf{Identifica las posiciones:} en tres bits, las posiciones valen 4, 2 y 1 de izquierda a derecha. Anótalas antes de operar.\\
Patrones & \textbf{Lee cada dígito:} un 1 enciende esa posición y suma su valor; un 0 la apaga y suma cero.\\
Abstracción & \textbf{Ignora lo que no cambia la respuesta:} el nombre del número no importa; solo importan los dígitos y sus posiciones.\\
Algoritmo de armado & \textbf{Suma en orden:} de izquierda a derecha, posición a posición, y escribe el resultado. Verifica que esté entre 0 y 7.\\
\end{tabularx}
\normalsize

\begin{drawbox}{milcMagenta}{El reto: ¿cuánto vale 110?}
\checkbox \textbf{Reto (Nivel B).} En binario, las tres posiciones valen 4, 2 y 1 (de izquierda a derecha). Cada posición está encendida (1) o apagada (0) y se suma lo encendido. Ejemplo: 101 = 4 + 0 + 1 = 5. ¿Cuánto vale \textbf{110}? \\[2mm]
\checkbox Marca: \quad \fbox{\phantom{x}}~3 \quad \fbox{\phantom{x}}~5 \quad \fbox{\phantom{x}}~6 \quad \fbox{\phantom{x}}~7 \\[2mm]
\checkbox Escribe \textbf{los pasos} que seguiste para llegar a la respuesta.
\end{drawbox}

\begin{softbox}{milcMagenta}{milcGris}{Plantilla de trabajo}
\textbf{Resuélvelo paso a paso.} Las posiciones de 110 valen: la primera (1) vale 4, la segunda (1) vale 2, la tercera (0) vale 1. Suma solo las encendidas: 4 + 2 + 0 = \textbf{6}. La respuesta correcta es 6. Pieza usada: \textbf{binario} --- leer el valor posicional y sumar los bits encendidos.
\end{softbox}

\begin{infoband}{milcMagenta}
\textbf{En tu cuaderno (Actividad 3 · APLICA).} Título: «Actividad 3 --- Mi reto de binario». \textbf{Formato:} respuesta marcada + 3 renglones de explicación. \textbf{Extensión:} la respuesta y los pasos seguidos, nombrando las posiciones que encendiste. \textbf{Sabes que terminaste cuando} tu explicación permite que alguien que no lo resolvió entienda cómo llegaste al 6.
\end{infoband}

\puente{Tienes dos evidencias. Ahora deja tu trabajo en el cuaderno con el formato que indica cada actividad, para mostrar que de verdad sabes codificar con dos estados.}

\begin{titlebox}{milcMostaza}
Producto de la sesión
\end{titlebox}
\textbf{Párrafo «Lo mío en dos estados» y tabla «Mi código de tambor» con 4 mensajes codificados}.
\begin{softbox}{milcMostaza}{milcGris}{Criterios de calidad}
(1) Resolviste el reto binario y marcaste la respuesta correcta (6 = 4 + 2 + 0). (2) Explicaste los pasos: nombraste el valor de cada posición encendida y los sumaste. (3) Tu párrafo «Lo mío en dos estados» conecta algo concreto de tu vida con la idea de codificar con 0 y 1. (4) Tu tabla «Mi código de tambor» tiene 4 mensajes distintos, cada uno con su código de golpes (G) y silencios (S), sin repetir ningún código. (5) Tu compañero pudo decodificar al menos un mensaje usando solo tu tabla, sin que se lo expliques.
\end{softbox}

\puente{Terminaste el producto. Detente a mirar qué cambió en ti --- no la nota, sino tu manera de ver los datos que te rodean.}

\begin{titlebox}{milcVino}
Fase 4 · Evaluación liberadora — cinco dimensiones
\end{titlebox}

\begin{softbox}{milcVino}{milcGris}{Sentido de la evaluación}
Evaluar no es solo revisar una nota. Es reconocer qué aprendí, qué cambió en mí, qué emociones manejé, cómo me relaciono con la ciudad y la comunidad, y qué le dejaré a quien viene detrás.
\end{softbox}

\textbf{Mi reflexión liberadora en cinco dimensiones.} Completa cada frase iniciada:

\small
\begin{tabularx}{\textwidth}{>{\bfseries\RaggedRight\arraybackslash}p{3.4cm}Y>{\RaggedRight\arraybackslash}p{4.6cm}}
\rowcolor{milcVino}\color{white}Dimensión & \color{white}\bfseries Mi respuesta (frame starter) & \color{white}\bfseries Algo que descubrí\\
1. Personal & Lo que más cambió en mí fue \linead{6.5cm} & \linead{4.2cm}\\
2. Emocional & La emoción más fuerte fue \linead{2.5cm} y la regulé \linead{3cm} & \linead{4.2cm}\\
3. Ciudadana & En democracia esto importa porque \linead{6cm} & \linead{4.2cm}\\
4. Local (Cartago) & En mi barrio sería útil para \linead{6.5cm} & \linead{4.2cm}\\
5. Intergeneracional & A mi abuela/abuelo le diría \linead{6.5cm} & \linead{4.2cm}\\
\end{tabularx}
\normalsize

\begin{infoband}{milcVino}
\textbf{Para la próxima sustentación.} Vuelve a esta tabla en seis meses. Verás que las respuestas cambian — no porque lo aprendido cambie, sino porque tú cambias. La evaluación liberadora es una conversación contigo a través del tiempo.
\end{infoband}


\puente{Cerremos pensando en grande: tres voces para entender qué significa, para ti y para tu comunidad, que toda tu vida digital quepa en dos estados.}

\begin{titlebox}{milcGradeDark}
Triángulo de pensamiento · cierre filosófico
\end{titlebox}

\begin{softbox}{milcVino}{milcGris}{Dussel · lente del nosotros}
\textit{``El saber del pueblo también es ciencia.''}

El tamborero afrocolombiano del Pacífico codificaba y transmitía mensajes con dos estados mucho antes de que existiera la palabra «dato». No llegó de afuera: tu región ya inventó informática, aunque no la llamara así. Reconocer eso es reconocer que tu comunidad produce conocimiento.

\textbf{¿Qué otras formas de representar información existen en tu comunidad --- canciones, tejidos, dibujos --- que nadie llama «datos» pero que funcionan igual?}
\end{softbox}
\linead{\linewidth}\\[1mm]
\linead{\linewidth}

\begin{softbox}{milcMostaza}{milcGris}{Estoicismo · Séneca · cuidado interior}
\textit{``No es pobre el que tiene poco, sino el que desea más. Lo esencial cabe en lo simple.''}

Con solo dos estados --- golpe y silencio, 0 y 1 --- se representa toda la información digital del planeta. No hace falta más: lo que importa es el acuerdo y el orden, no la cantidad de símbolos. Lo complejo, bien partido, cabe en lo simple.

\textbf{¿Cuántas veces complico de más algo que, bien pensado, se reduce a unos pocos elementos esenciales bien ordenados?}
\end{softbox}
\linead{\linewidth}\\[1mm]
\linead{\linewidth}

\begin{softbox}{milcTurquesa}{milcGris}{Floridi · lente de la infoesfera}
\textit{``Todo lo que habita la infoesfera es, en el fondo, información codificada.''}

Tu foto, tu voz, tu nombre: todo se convierte en patrones de 0 y 1 que alguien guarda en un servidor. Entender cómo se representan los datos es el primer paso para decidir qué compartes y con quién.

\textbf{¿Qué tanto de mi vida ya está convertido en datos --- en unos y ceros --- sin que yo lo haya decidido conscientemente?}
\end{softbox}
\linead{\linewidth}\\[1mm]
\linead{\linewidth}

\begin{titlebox}{milcVino}
Compromiso final
\end{titlebox}

\small
\begin{tabularx}{\textwidth}{>{\RaggedRight\arraybackslash}p{3.0cm}YYY}
\rowcolor{milcVino}\color{white}\bfseries Criterio & \color{white}\bfseries Lo logré & \color{white}\bfseries En proceso & \color{white}\bfseries Necesito apoyo\\
Comprensión del tema & Explico ideas centrales con mis palabras. & Reconozco ideas pero falta precisión. & Necesito repasar conceptos base.\\
Pensamiento computacional & Apliqué los 4 pilares al diseñar mi producto. & Apliqué algunos pilares. & Necesito releer la fase 2.\\
Aplicación y producto & Mi producto cumple los criterios y se puede revisar. & Producto funcional con ajustes pendientes. & Aún en construcción.\\
Reflexión 5 dimensiones & Respondí cada dimensión con honestidad. & Respondí algunas con detalle. & Mis respuestas son muy generales.\\
Triángulo de pensamiento & Conecté las 3 voces con mi trabajo real. & Conecté al menos una. & No relaciono las voces con mi proceso.\\
\end{tabularx}
\normalsize

\textbf{Hoy entendí que:}\\[1mm]
\linead{\linewidth}\\[1mm]
\linead{\linewidth}

\textbf{Mi compromiso para la próxima sesión es:}\\[1mm]
\linead{\linewidth}\\[1mm]
\linead{\linewidth}

\begin{softbox}{milcMostaza}{milcGris}{Cita de despedida · Marco Aurelio}
\textit{``El obstáculo en el camino se vuelve el camino. Lo que estorba al camino se vuelve camino.''}\\[1mm]
Cada error de hoy, bien observado, se convierte en pieza del camino que pisarás mañana.
\end{softbox}

\small
\begin{tabularx}{\textwidth}{>{\bfseries}p{3.6cm}Y}
\rowcolor{milcGradeDark!12}Nombre completo: & \linead{10cm}\\
Fecha: & \linead{4cm} \quad Curso/grupo: \linead{4cm}\\
Mi firma: & \linead{9cm}\\
\end{tabularx}
\normalsize

\begin{infoband}{milcGradeDark}
\textbf{Para la próxima sesión.} Busca en tu casa o en tu barrio un ejemplo de cada representación: algo que vaya en orden (lista), algo con filas y columnas (tabla) y algo que se ramifique (árbol). Trae también un código de dos estados que hayas inventado tú --- no tiene que ser de tambor: puede ser con luces, golpes en la mesa o colores. Explica ante el grupo qué significa cada combinación.
\end{infoband}

\end{document}
